\section{The next discriminating comparison}
\label{sec:next-study}

The primary question is whether model-corrected mixture proposals and
ancestor averaging reduce model-score error at equal compute. Proposal
quality, score construction, and mixture-gradient variance must be separated.
Otherwise a successful run could be attributed to the wrong mechanism, just
as a passing derivative check previously obscured a changed scalar.

The smallest useful sequence begins with three calculations. First, compare
\eqref{eq:model-marginal-weights} and
\eqref{eq:fully-adapted-mixture} with direct Gaussian integration on a fixed
incoming cloud. Second, compare the complete-data and backward-recursion
scores with exact enumeration on a small finite-state model, including a
parameter-dependent initial law. Third, compare IWSG, centered IWSG, and
\eqref{eq:hybrid-estimator} for the same declared mixture expectation.
These calculations test identities and failure mechanisms; their particle
counts need only be large enough to exercise the operation being tested.
They are not score-quality campaigns or reasons to select a production
particle count.

The subsequent accuracy study should use a linear-Gaussian model with an
independent Kalman score, a short Gaussian-mixture model whose modes can be
enumerated exactly, and a scalar nonlinear model with a numerically converged
quadrature score. Kalman checks the basic integration and complete derivatives.
The mixture model tests the separated modes that motivate Younis's method.
The nonlinear model tests how proposal approximation affects the result.
A one-step or unimodal success cannot establish the latter two properties.
Models with unresolved initialization or support identity must be repaired
before their score results are interpreted.

For each model, change one mechanism at a time: ordinary ancestor sampling
versus marginal mixture weights; a genealogical complete-data score versus
the full ancestor average; and, where an expectation derivative is actually
used, all-IWSG versus the selective estimator. Include the enhanced
combination only after the component comparisons. The exact full-mixture
reference remains necessary when evaluating a later sparse or sampled
approximation.

The cheap alternatives are constructed from the model before claim data are
examined. A bootstrap PF proposes from the physical transition and uses a
proper complete-data score estimator. A locally adapted PF uses the current
observation to guide each ancestor's proposal and keeps the exact importance
weight. In nearly linear regimes, EKF and UKF approximate-likelihood scores
test whether inexpensive Gaussian propagation already suffices; their
derivatives must include the full declared Gaussian approximation. A
scope-tuned canonical LEDH filter supplies the project comparator. The exact
Kalman, mixture-enumeration, or quadrature solution supplies the oracle.
Neither a weak fixed-ancestry derivative nor an untuned numerical default
should stand in for all conventional particle-score methods.

Evaluate these comparisons conditionally. Weak and strong observations probe
proposal concentration; overlapping and separated modes probe mixture
gradients; low and high persistence probe loss of ancestor diversity; short
and long horizons probe accumulation; near-zero and large score components
probe the error scale. These situations must be defined from the model and
research question before using validation data. Failure against a cheap
alternative in a salient situation rejects promotion under the repository's
sanity-check rule, even if the method wins an unconditional average. The
adversaries are falsification checks, not extra tuning targets.

For dataset $d$ and independent filter randomization $r$, denote the estimate
by $\widehat s_{d,r}$ and the oracle by $s_d$. Replication within each fixed
dataset estimates $b_d$ and $V_d$ from
\eqref{eq:conditional-score-error}; the distribution over datasets measures
generalization. Pair methods within datasets and respect this nesting in
confidence intervals. A path bootstrap with only one filter realization per
path cannot estimate each dataset's conditional bias. Freeze parameter
coordinate scales before comparison, because a relative error divided by a
nearly zero oracle score is unstable. Change $N$ and $T$ independently, and
report both equal-particle and equal-measured-compute comparisons.

Before a serious run, choose a practically meaningful reduction $\delta$ and
test a predeclared upper confidence bound for
\begin{equation}
 \Delta=\mathbb E\!\left[
    \|\widehat s_{\rm candidate}-s\|_D^2
    -(1-\delta)\|\widehat s_{\rm comparator}-s\|_D^2\right],
 \qquad \|v\|_D^2=v^\mathsf T Dv,
 \label{eq:next-study-criterion}
\end{equation}
where $D$ is the fixed coordinate scaling and the expectation follows the
specified sampling of datasets and filter randomizations. The upper bound must be below
zero, subject to validity and conditional-comparison vetoes. The old ten
percent criterion is not automatically a justified threshold for every new
model. Its practical meaning, replication requirement, multiplicity handling,
and total compute budget must be stated for the new study. A power pilot that
requires more replications than the budget permits leads to an inconclusive
study, not a forced ranking among viable candidates.

All material numerical choices need a provenance and failure explanation.
The physical transition covariance comes from the model; proposal bandwidth,
defensive-mixture fraction, covariance marks, chart, regularization, and
ancestor-sampling precision are numerical hypotheses. The earliest checks
should expose wrong support, missing proposal mass, large importance ratios,
or a gradient dominated by inverse bandwidth. Calibration and validation
must be disjoint for each bound model/horizon/particle/backend scope. The
initial law is part of that scope. Retaining an old setting for comparability
does not make it a tuned default, and a failed holdout cannot become the next
calibration set.

Numerical validity is assessed before score ranking. Nonfinite values,
incorrect density ratios, support violations, missing initial-law terms,
an invalid reference, or an omitted sampling-law derivative invalidate the
affected comparison. ESS, mode overlap, cap activity, tangent norms, bandwidth,
memory, and runtime explain its behavior; they do not replace the model-score
criterion. A valid candidate that loses is a candidate rejection. A later
prespecified repair may continue unless the target, evidence, or total budget
has become invalid.

BayesFilter implementation remains TensorFlow/TFP with analytical recursive
scores. Autodiff and finite differences are independent diagnostics. The
production direction remains GPU/XLA with the declared memory-growth and
transport-chunk policies; the earlier float64 and no-TF32 tests are reference
exceptions. New mixture evaluation must expose its pair count and memory
behavior. Batch wrappers must preserve the mathematical calculation and,
for NeuTra training, actually evaluate a batch rather than apply a scalar
target independently to each row. None of the proposed new calculations is
promoted by this document.

The most useful next evidence is therefore modest but decisive: correct
model-density weights and complete initialization, an independently checked
score recursion, and a powered comparison against simple alternatives. The
strongest alternative explanation for any improvement is better proposal
adaptation or more effective computation, rather than the gradient identity.
The component comparisons distinguish those explanations. Loss to a tuned
conventional score estimator at equal compute would overturn the practical
priority of the new combination without invalidating its algebra.
